% META: {"id": "TSPE-DERCONV-EX-030", "chapitre": "TSPE-DERIVATION-CONVEXITE", "type_objet": "exercice", "capacites": ["TSPE-DERCONV-C4"], "capacites_codes": ["C4"], "methodes": ["M4"], "parcours": 2, "competences": ["raisonner","calculer"], "duree_min": 15, "mode_creation": "ex_nihilo", "sources_inspiration": [], "fichier_tex": "chapitres/TSPE-DERIVATION-CONVEXITE/exercices/TSPE-DERCONV-EX-030.tex", "corrige_tex": "chapitres/TSPE-DERIVATION-CONVEXITE/corriges/TSPE-DERCONV-CO-030.tex", "status": "generated"}
% BEGIN-VERIFY
% from sympy import *
% x = symbols('x')
% # f' donnee par son tableau: croissante sur R-, decroissante sur R+, f'(0)=2
% # cela suggere f'(x) = 2 - x^2, donc f(x) = 2x - x^3/3
% fp = 2 - x**2
% f = integrate(fp, x)
% assert f == 2*x - x**3/3
% fpp = diff(f, x, 2)
% assert fpp == -2*x
% assert solve(fp, x) == [-sqrt(2), sqrt(2)]
% assert solve(fpp, x) == [0]
% END-VERIFY
\begin{exercice}{TSPE-DERCONV-EX-030}{2}{15}
On donne le tableau de variations de $f'$ : $f'$ est croissante sur $\left]-\infty, 0\right]$, decroissante sur $\left[0, +\infty\right[$, avec $f'(0) = 2$ et $f'(\pm\sqrt{2}) = 0$. On suppose $f'(x) = 2 - x^2$.
\begin{enumerate}
  \item En deduire une expression possible de $f$.
  \item Determiner les intervalles de convexite de $f$.
  \item Esquisser la courbe de $f$.
\end{enumerate}

\end{exercice}
